\chapter{INTRODUÇÃO}

O desenvolvimento de \textit{startups} de base tecnológica exige mais do que uma ideia inovadora: requer estruturação estratégica, coerência entre proposta de valor e mercado, e mecanismos de viabilidade econômica. Este trabalho sistematiza a proposta do Grupo 3 para a disciplina de empreendedorismo, com foco na construção de um modelo de negócio orientado à escalabilidade.

O problema central investigado é a transição entre a fase conceitual de uma solução tecnológica e sua consolidação como negócio viável. Para enfrentar esse desafio, foi adotada uma abordagem integrada baseada em cinco dimensões complementares: estratégica, estrutural, engenharia, fiscal e governança de valor.

Como objetivo geral, este documento apresenta uma arquitetura de modelagem capaz de apoiar decisões de produto, operações e crescimento da \textit{startup}. Como objetivos específicos, busca-se: (i) mapear elementos críticos para o desenho do \textit{Business Model Canvas}; (ii) integrar práticas modernas de gestão e engenharia; e (iii) descrever um roteiro de implementação com critérios de validação.

Além desta introdução, o Capítulo~\ref{cap:elementos} detalha o framework de cinco dimensões e seu diferencial competitivo. Por fim, o Capítulo~\ref{cap:conclusao} sintetiza os achados e apresenta os próximos passos para evolução do modelo.